% Paper E: Fully Fused Logic Accelerator - FPGA Implementation
% The "Hardware Paper" - Mapping constraint propagation to FPGA primitives
%
% For God so loved the world that he gave his only begotten Son,
% that whoever believes in him should not perish but have eternal life.
% John 3:16

\documentclass[11pt,a4paper]{article}

\usepackage[utf8]{inputenc}
\usepackage{amsmath,amssymb,amsthm}
\usepackage{booktabs}
\usepackage{hyperref}
\usepackage{xspace}
\usepackage{listings}
\usepackage{xcolor}
\usepackage{tikz}
\usepackage{pgfplots}
\usepackage{float}
\pgfplotsset{compat=1.18}
\usetikzlibrary{shapes,arrows,positioning,fit,calc}

% Theorem environments
\newtheorem{theorem}{Theorem}
\newtheorem{lemma}[theorem]{Lemma}
\newtheorem{definition}[theorem]{Definition}
\newtheorem{proposition}[theorem]{Proposition}

% Code listing style for Calyx/Clash
\lstdefinestyle{calyxstyle_chirho}{
  basicstyle=\ttfamily\small,
  keywordstyle=\color{blue}\bfseries,
  commentstyle=\color{gray},
  stringstyle=\color{orange},
  numbers=left,
  numberstyle=\tiny\color{gray},
  breaklines=true,
  frame=single,
  captionpos=b
}

% Include shared definitions
% Shared Definitions for miniKanren 1-Bit Paper Suite
% Include with: % Shared Definitions for miniKanren 1-Bit Paper Suite
% Include with: % Shared Definitions for miniKanren 1-Bit Paper Suite
% Include with: \input{../shared_chirho/shared_defs_chirho}
%
% For God so loved the world that he gave his only begotten Son,
% that whoever believes in him should not perish but have eternal life.
% John 3:16

% ============================================================================
% Core Definitions (shared across all papers)
% ============================================================================

% The Finite-Domain miniKanren Kernel
\newcommand{\fdmk}{FD-miniKanren\xspace}  % Finite-Domain miniKanren

% ============================================================================
% Scope Clarification (include in each paper's introduction)
% ============================================================================

\newcommand{\scopeclarification}{%
\textbf{Scope Clarification.}
This paper addresses the \emph{Finite-Domain miniKanren Kernel} (\fdmk),
where variable domains are finite bitmasks.
Standard miniKanren unification over unbounded term structures uses a separate
reference implementation; the bit-parallel approach accelerates the constraint
propagation core, not full structural unification.
}

% ============================================================================
% Common Notation
% ============================================================================

% Domains
\newcommand{\dom}[1]{D_{#1}}          % Domain of variable x
\newcommand{\domand}{\land}            % Domain intersection (AND)
\newcommand{\domor}{\lor}              % Domain union (OR)

% Search
\newcommand{\state}{(\vec{D}, \sigma)} % Search state
\newcommand{\findChirho}{\text{find}}      % Union-find find
\newcommand{\unionChirho}{\text{union}}    % Union-find union

% Tensors
\newcommand{\tensor}[1]{T_{#1}}        % Tensor for relation R
\newcommand{\contract}{\circledast}    % Tensor contraction

% Semirings
\newcommand{\softand}{\odot}           % Soft AND (multiply)
\newcommand{\softor}{\oplus}           % Soft OR (prob sum)

% ============================================================================
% Hierarchical Domain Scaling
% ============================================================================

% Domain type naming (semantic macros for consistency across papers)
% Using semantic names to avoid confusion with Rust type identifiers
\newcommand{\domainSmallChirho}{\texttt{BitVec64Chirho}}
\newcommand{\domainMediumChirho}{\texttt{Hierarchical4kChirho}}
\newcommand{\domainLargeChirho}{\texttt{Hierarchical256kChirho}}
\newcommand{\domainDiffChirho}{\texttt{DiffHierarchical4kChirho}}

% Measured performance values (from Rust benchmarks, Jan 2026)
\newcommand{\domainSmallTimeChirho}{420 ps}
\newcommand{\domainMediumTimeChirho}{39 ns}
\newcommand{\domainLargeTimeChirho}{2.5 $\mu$s}
\newcommand{\domainLargeOpsChirho}{400K ops/sec}

% ============================================================================
% Paper Suite Reference
% ============================================================================

% Use these to cite sibling papers
\newcommand{\paperA}{Paper A: Bit-Parallel Finite-Domain Relational Search}
\newcommand{\paperB}{Paper B: Relations as Sparse Boolean Tensors}
\newcommand{\paperC}{Paper C: Bridging Infinite Domains with Hash-Consed Term IDs}
\newcommand{\paperD}{Paper D: Tabling and Mode-Driven Goal Scheduling}
\newcommand{\paperE}{Paper E: Fully Fused Logic Accelerator}
\newcommand{\paperF}{Paper F: Differentiable Constraint Solving via Semiring Relaxation}

% ============================================================================
% Acknowledgments Template
% ============================================================================

\newcommand{\standardacknowledgments}{%
Above all, I thank my Lord and Savior Jesus Christ for saving a wretch like me.

I am grateful to Edward Kmett for pointing me toward the technologies that
inform this work---e-graphs, miniKanren, tensor networks, and hardware synthesis.

I also thank my father, Roger Cowan, for his patient support.

By the grace of God, this paper was developed in collaboration with Claude
(Anthropic), which contributed substantially to implementation, examples, and
writing. Google Gemini and OpenAI's GPT-5.2 Codex provided valuable feedback
on drafts.
}

% ============================================================================
% Example environment (for papers that use it)
% ============================================================================
\newenvironment{exampleChirho}{\par\noindent\textbf{Example.}\itshape}{\par}

% Soli Deo Gloria

%
% For God so loved the world that he gave his only begotten Son,
% that whoever believes in him should not perish but have eternal life.
% John 3:16

% ============================================================================
% Core Definitions (shared across all papers)
% ============================================================================

% The Finite-Domain miniKanren Kernel
\newcommand{\fdmk}{FD-miniKanren\xspace}  % Finite-Domain miniKanren

% ============================================================================
% Scope Clarification (include in each paper's introduction)
% ============================================================================

\newcommand{\scopeclarification}{%
\textbf{Scope Clarification.}
This paper addresses the \emph{Finite-Domain miniKanren Kernel} (\fdmk),
where variable domains are finite bitmasks.
Standard miniKanren unification over unbounded term structures uses a separate
reference implementation; the bit-parallel approach accelerates the constraint
propagation core, not full structural unification.
}

% ============================================================================
% Common Notation
% ============================================================================

% Domains
\newcommand{\dom}[1]{D_{#1}}          % Domain of variable x
\newcommand{\domand}{\land}            % Domain intersection (AND)
\newcommand{\domor}{\lor}              % Domain union (OR)

% Search
\newcommand{\state}{(\vec{D}, \sigma)} % Search state
\newcommand{\findChirho}{\text{find}}      % Union-find find
\newcommand{\unionChirho}{\text{union}}    % Union-find union

% Tensors
\newcommand{\tensor}[1]{T_{#1}}        % Tensor for relation R
\newcommand{\contract}{\circledast}    % Tensor contraction

% Semirings
\newcommand{\softand}{\odot}           % Soft AND (multiply)
\newcommand{\softor}{\oplus}           % Soft OR (prob sum)

% ============================================================================
% Hierarchical Domain Scaling
% ============================================================================

% Domain type naming (semantic macros for consistency across papers)
% Using semantic names to avoid confusion with Rust type identifiers
\newcommand{\domainSmallChirho}{\texttt{BitVec64Chirho}}
\newcommand{\domainMediumChirho}{\texttt{Hierarchical4kChirho}}
\newcommand{\domainLargeChirho}{\texttt{Hierarchical256kChirho}}
\newcommand{\domainDiffChirho}{\texttt{DiffHierarchical4kChirho}}

% Measured performance values (from Rust benchmarks, Jan 2026)
\newcommand{\domainSmallTimeChirho}{420 ps}
\newcommand{\domainMediumTimeChirho}{39 ns}
\newcommand{\domainLargeTimeChirho}{2.5 $\mu$s}
\newcommand{\domainLargeOpsChirho}{400K ops/sec}

% ============================================================================
% Paper Suite Reference
% ============================================================================

% Use these to cite sibling papers
\newcommand{\paperA}{Paper A: Bit-Parallel Finite-Domain Relational Search}
\newcommand{\paperB}{Paper B: Relations as Sparse Boolean Tensors}
\newcommand{\paperC}{Paper C: Bridging Infinite Domains with Hash-Consed Term IDs}
\newcommand{\paperD}{Paper D: Tabling and Mode-Driven Goal Scheduling}
\newcommand{\paperE}{Paper E: Fully Fused Logic Accelerator}
\newcommand{\paperF}{Paper F: Differentiable Constraint Solving via Semiring Relaxation}

% ============================================================================
% Acknowledgments Template
% ============================================================================

\newcommand{\standardacknowledgments}{%
Above all, I thank my Lord and Savior Jesus Christ for saving a wretch like me.

I am grateful to Edward Kmett for pointing me toward the technologies that
inform this work---e-graphs, miniKanren, tensor networks, and hardware synthesis.

I also thank my father, Roger Cowan, for his patient support.

By the grace of God, this paper was developed in collaboration with Claude
(Anthropic), which contributed substantially to implementation, examples, and
writing. Google Gemini and OpenAI's GPT-5.2 Codex provided valuable feedback
on drafts.
}

% ============================================================================
% Example environment (for papers that use it)
% ============================================================================
\newenvironment{exampleChirho}{\par\noindent\textbf{Example.}\itshape}{\par}

% Soli Deo Gloria

%
% For God so loved the world that he gave his only begotten Son,
% that whoever believes in him should not perish but have eternal life.
% John 3:16

% ============================================================================
% Core Definitions (shared across all papers)
% ============================================================================

% The Finite-Domain miniKanren Kernel
\newcommand{\fdmk}{FD-miniKanren\xspace}  % Finite-Domain miniKanren

% ============================================================================
% Scope Clarification (include in each paper's introduction)
% ============================================================================

\newcommand{\scopeclarification}{%
\textbf{Scope Clarification.}
This paper addresses the \emph{Finite-Domain miniKanren Kernel} (\fdmk),
where variable domains are finite bitmasks.
Standard miniKanren unification over unbounded term structures uses a separate
reference implementation; the bit-parallel approach accelerates the constraint
propagation core, not full structural unification.
}

% ============================================================================
% Common Notation
% ============================================================================

% Domains
\newcommand{\dom}[1]{D_{#1}}          % Domain of variable x
\newcommand{\domand}{\land}            % Domain intersection (AND)
\newcommand{\domor}{\lor}              % Domain union (OR)

% Search
\newcommand{\state}{(\vec{D}, \sigma)} % Search state
\newcommand{\findChirho}{\text{find}}      % Union-find find
\newcommand{\unionChirho}{\text{union}}    % Union-find union

% Tensors
\newcommand{\tensor}[1]{T_{#1}}        % Tensor for relation R
\newcommand{\contract}{\circledast}    % Tensor contraction

% Semirings
\newcommand{\softand}{\odot}           % Soft AND (multiply)
\newcommand{\softor}{\oplus}           % Soft OR (prob sum)

% ============================================================================
% Hierarchical Domain Scaling
% ============================================================================

% Domain type naming (semantic macros for consistency across papers)
% Using semantic names to avoid confusion with Rust type identifiers
\newcommand{\domainSmallChirho}{\texttt{BitVec64Chirho}}
\newcommand{\domainMediumChirho}{\texttt{Hierarchical4kChirho}}
\newcommand{\domainLargeChirho}{\texttt{Hierarchical256kChirho}}
\newcommand{\domainDiffChirho}{\texttt{DiffHierarchical4kChirho}}

% Measured performance values (from Rust benchmarks, Jan 2026)
\newcommand{\domainSmallTimeChirho}{420 ps}
\newcommand{\domainMediumTimeChirho}{39 ns}
\newcommand{\domainLargeTimeChirho}{2.5 $\mu$s}
\newcommand{\domainLargeOpsChirho}{400K ops/sec}

% ============================================================================
% Paper Suite Reference
% ============================================================================

% Use these to cite sibling papers
\newcommand{\paperA}{Paper A: Bit-Parallel Finite-Domain Relational Search}
\newcommand{\paperB}{Paper B: Relations as Sparse Boolean Tensors}
\newcommand{\paperC}{Paper C: Bridging Infinite Domains with Hash-Consed Term IDs}
\newcommand{\paperD}{Paper D: Tabling and Mode-Driven Goal Scheduling}
\newcommand{\paperE}{Paper E: Fully Fused Logic Accelerator}
\newcommand{\paperF}{Paper F: Differentiable Constraint Solving via Semiring Relaxation}

% ============================================================================
% Acknowledgments Template
% ============================================================================

\newcommand{\standardacknowledgments}{%
Above all, I thank my Lord and Savior Jesus Christ for saving a wretch like me.

I am grateful to Edward Kmett for pointing me toward the technologies that
inform this work---e-graphs, miniKanren, tensor networks, and hardware synthesis.

I also thank my father, Roger Cowan, for his patient support.

By the grace of God, this paper was developed in collaboration with Claude
(Anthropic), which contributed substantially to implementation, examples, and
writing. Google Gemini and OpenAI's GPT-5.2 Codex provided valuable feedback
on drafts.
}

% ============================================================================
% Example environment (for papers that use it)
% ============================================================================
\newenvironment{exampleChirho}{\par\noindent\textbf{Example.}\itshape}{\par}

% Soli Deo Gloria


\title{Fully Fused Logic Accelerator:\\
FPGA Implementation of Bit-Parallel Constraint Propagation\\
\large The Hardware Paper}

\author{
  L.J. Brian Cowan \\
  \texttt{loveJesus@loveJes.us}
}

\date{Draft - \today}

\begin{document}

\maketitle

\begin{abstract}
We demonstrate that the bit-parallel constraint propagation kernel of
\fdmk{} maps directly to FPGA primitives: lookup tables (LUTs) implement
domain intersection, content-addressable memory (CAM) enables hash-consed
term lookup, and registers store search state. The result is a fully fused
logic accelerator where unification executes in a single clock cycle.

We present: (1) an architecture diagram showing the mapping from miniKanren
operations to hardware components; (2) a resource model predicting LUT,
register, and block RAM consumption; (3) simulation results from Calyx IR
compiled to Verilog and verified in Verilator (target: Xilinx Artix-7); and
(4) a discussion of limitations, including the current software-only
implementation of tabling and mode analysis.

\textbf{Headline Result.}
The Clash-generated searchEngineChirho achieves \textbf{40.5 million
unifications per second} in batch mode with \textbf{24.7~ns per-operation
latency}---a \textbf{6,679$\times$ speedup} over CPU for batch workloads.
Synthesis (Yosys 0.61) reports 43,136 cells fitting Cyclone~V at $\sim$8\%
utilization.

\textbf{Production-Scale Training Benchmarks (January 2026).}
We validated production-scale training workloads on AWS F2 with HBM batching:
\textbf{3.3 million$\times$} speedup for 1000-variable SAT training (7.9 sec CPU
$\to$ 2.4 $\mu$s FPGA) and \textbf{1.5 million$\times$} for 256-dimension
attention training. These results use Q16.16 fixed-point arithmetic with
Gumbel-softmax reparameterization entirely in FPGA HBM.

\textbf{Physical Validation (January 2026).}
We successfully synthesized and deployed the design on AWS F2 instances
(AMD VU47P-HBM FPGA) using Vivado 2025.1. The design achieves \textbf{timing
closure at 250~MHz} with +0.159~ns slack. Two Amazon FPGA Images were created:
\texttt{agfi-04abde24231f6775c} (initial validation) and
\texttt{agfi-05988b0b1980d6d2f} (comprehensive benchmarks).

\textbf{Verified Benchmarks.}
We ran verified benchmarks on real data (Greek New Testament, 137,498 words)
with CPU comparison. Key results: TestForge (26,288$\times$ speedup for 10K
records), ConfigGuard (9,657$\times$ for 5K files), and Philologos Greek NT
searches (6,679$\times$ for 8000-search batch).
\end{abstract}

% ============================================================================
\section{Introduction}
\label{sec:intro_chirho}
% ============================================================================

\scopeclarification

Logic programming systems like miniKanren and Prolog have traditionally been
implemented in software, with performance limited by memory allocation, pointer
chasing, and sequential execution. Our earlier work (Papers A--D) established
that finite-domain constraint propagation can be reduced to bit-parallel
operations: domain intersection is bitwise AND, failure detection is zero-test,
and branching is bit extraction.

This paper asks: \emph{Can these operations be mapped directly to FPGA
primitives?} We answer affirmatively by presenting a hardware architecture
that executes the core constraint propagation loop with:

\begin{itemize}
  \item \textbf{Single-cycle unification}: 64-bit AND via 6-input LUTs
  \item \textbf{O(1) term lookup}: CAM-based hash consing for structural sharing
  \item \textbf{Parallel branching}: Simultaneous exploration of multiple search paths
  \item \textbf{Deterministic timing}: No garbage collection or cache misses
\end{itemize}

The architecture is implemented in two complementary hardware description
approaches:
\begin{itemize}
  \item \textbf{Calyx IR}: A compiler infrastructure for hardware accelerators,
    allowing high-level specification with guaranteed correct lowering to Verilog
  \item \textbf{Clash}: A functional HDL based on Haskell, providing strong
    type guarantees and high-level abstractions while generating synthesizable RTL
\end{itemize}

% ============================================================================
\section{Background: FPGA Primitives}
\label{sec:background_chirho}
% ============================================================================

Modern FPGAs (e.g., Xilinx 7-series, Intel Cyclone V) provide three key
primitive types relevant to our accelerator:

\subsection{Lookup Tables (LUTs)}
\label{subsec:luts_chirho}

A $k$-input LUT implements any Boolean function of $k$ variables in O(1) time
by storing a $2^k$-entry truth table. For 6-input LUTs (standard in modern
FPGAs), 64-bit operations require $\lceil 64/1 \rceil = 64$ LUTs operating in
parallel.

\textbf{Key insight}: Domain intersection ($D_x \land D_y$) maps to 64
parallel 2-input AND operations, each consuming a single LUT.

\subsection{Block RAM (BRAM)}
\label{subsec:bram_chirho}

Block RAM provides high-bandwidth, dual-port memory. A typical Artix-7 BRAM
block offers 36 Kbit capacity, configurable as 1K$\times$36 or 2K$\times$18.

\textbf{Application}: Term storage for hash consing. With 34-bit terms
(2-bit tag + 16-bit left + 16-bit right), a single BRAM stores 1K terms.

\subsection{Content-Addressable Memory (CAM)}
\label{subsec:cam_chirho}

CAM enables O(1) lookup by content rather than address. While FPGAs lack
native CAM, it can be emulated using hash tables with LUT-based collision
handling.

\textbf{Application}: Term deduplication in hash consing. A 64-entry CAM
with 6-bit hash provides $\approx$93\% hit rate for typical workloads.

% ============================================================================
\section{Architecture}
\label{sec:architecture_chirho}
% ============================================================================

Figure~\ref{fig:arch_chirho} presents the high-level architecture of the
logic accelerator.

\begin{figure}[ht]
\centering
\begin{tikzpicture}[
  block/.style={rectangle, draw, minimum width=2.5cm, minimum height=1cm, align=center},
  arrow/.style={->, >=stealth, thick}
]

% Components
\node[block] (cmd) at (0, 4) {Command\\Interface};
\node[block] (domains) at (-3, 2) {Domain\\Registers\\(64-bit $\times$ N)};
\node[block] (unify) at (0, 2) {Unify Unit\\(AND + Zero)};
\node[block] (branch) at (3, 2) {Branch Unit\\(Fork + Stack)};
\node[block] (hashcons) at (0, 0) {Hash Cons\\(CAM + BRAM)};
\node[block] (resp) at (0, -2) {Response\\Interface};

% Arrows
\draw[arrow] (cmd) -- (domains);
\draw[arrow] (cmd) -- (unify);
\draw[arrow] (cmd) -- (branch);
\draw[arrow] (domains) -- (unify);
\draw[arrow] (unify) -- (domains);
\draw[arrow] (unify) -- (branch);
\draw[arrow] (branch) -- (domains);
\draw[arrow] (unify) -- (hashcons);
\draw[arrow] (hashcons) -- (unify);
\draw[arrow] (domains) -- (resp);
\draw[arrow] (unify) -- (resp);

\end{tikzpicture}
\caption{High-level architecture of the logic accelerator. Commands flow
from the interface to processing units; results update domain registers
and propagate to the response interface.}
\label{fig:arch_chirho}
\end{figure}

\subsection{Domain Registers}
\label{subsec:domain_regs_chirho}

Each logic variable is assigned a 64-bit register storing its domain bitmask.
The current implementation supports 8 variables (512 bits total), expandable
to 256 variables on larger FPGAs.

\subsection{Unification Unit}
\label{subsec:unify_unit_chirho}

The core computation: given two 64-bit domains $D_1$ and $D_2$, compute
$D' = D_1 \land D_2$ and check $D' \neq 0$. This requires:

\begin{itemize}
  \item 64 LUTs for the AND operation (parallel)
  \item 6 LUTs for the 64-input OR reduction (failure check)
  \item 1 register for the result
\end{itemize}

Total latency: 1 clock cycle. See Listing~\ref{lst:unify_chirho}.

\begin{lstlisting}[style=calyxstyle_chirho, caption={Unification unit in Calyx IR (excerpt from \texttt{calyx\_chirho/domain\_chirho.futil})}, label={lst:unify_chirho}]
component unify_unit_chirho(
  @go go: 1,
  d1_chirho: 64,
  d2_chirho: 64
) -> (
  @done done: 1,
  result_chirho: 64,
  failed_chirho: 1
) {
  cells {
    result_reg_chirho = std_reg(64);
    and_gate_chirho = std_and(64);
    zero_chirho = std_const(64, 0);
    eq_chirho = std_eq(64);
    failed_reg_chirho = std_reg(1);
  }
  // ... control: seq { compute_and; check_failure; }
}
\end{lstlisting}

\subsection{Branch Unit}
\label{subsec:branch_unit_chirho}

For non-deterministic search, the branch unit splits a domain into its lowest
set bit and remainder:

\begin{align}
  \text{lowest} &= x \land (-x) \\
  \text{rest} &= x \land (x - 1)
\end{align}

A 16-entry stack stores alternative branches for backtracking. Stack push/pop
executes in 1 cycle.

\subsection{Hash Consing Unit}
\label{subsec:hashcons_unit_chirho}

The hash consing unit provides structural sharing for terms, essential when
bridging to infinite domains via demand-driven term creation:

\begin{itemize}
  \item \textbf{Term format}: 34 bits (2-bit tag + 16-bit left + 16-bit right)
  \item \textbf{CAM}: 64 entries, 6-bit hash (XOR-based)
  \item \textbf{BRAM}: 1024 term slots
  \item \textbf{Latency}: 1 cycle for lookup, 2 cycles for insert
\end{itemize}

Listing~\ref{lst:hashcons_chirho} shows the Clash implementation.

\begin{lstlisting}[style=calyxstyle_chirho, caption={Hash consing types in Clash (excerpt from \texttt{clash\_chirho/HashConsChirho.hs})}, label={lst:hashcons_chirho}]
data TermChirho = TermChirho
  { tagChirho   :: TermTagChirho
  , leftChirho  :: TermIdChirho
  , rightChirho :: TermIdChirho
  } deriving (Generic, NFDataX)

hashTermChirho :: TermChirho -> Unsigned 6
hashTermChirho (TermChirho t l r) =
  let tagBits = resize (unpack (pack t))
  in truncateB $ tagBits `xor` l `xor` r
\end{lstlisting}

% ============================================================================
\section{Resource Model}
\label{sec:resource_chirho}
% ============================================================================

Table~\ref{tab:resources_chirho} summarizes the projected resource consumption
for an 8-variable configuration targeting Xilinx Artix-7 (XC7A35T).

\begin{table}[ht]
\centering
\caption{Projected FPGA resource consumption (8 variables, 64-bit domains)}
\label{tab:resources_chirho}
\begin{tabular}{lrrr}
\toprule
\textbf{Component} & \textbf{LUTs} & \textbf{Registers} & \textbf{BRAM} \\
\midrule
Domain Registers (8$\times$64) & 0 & 512 & 0 \\
Unification Unit & 70 & 65 & 0 \\
Branch Unit & 140 & 130 & 0 \\
Backtrack Stack (16 entries) & 64 & 8,320 & 0 \\
Hash Cons CAM (64 entries) & 2,048 & 2,240 & 0 \\
Hash Cons BRAM (1K terms) & 0 & 0 & 2 \\
Control Logic & 200 & 50 & 0 \\
\midrule
\textbf{Total} & \textbf{2,522} & \textbf{11,317} & \textbf{2} \\
\midrule
XC7A35T Available & 20,800 & 41,600 & 50 \\
\textbf{Utilization} & \textbf{12.1\%} & \textbf{27.2\%} & \textbf{4.0\%} \\
\bottomrule
\end{tabular}
\end{table}

\textbf{Variable Scaling}: Doubling to 16 variables approximately doubles register
count but has minimal impact on LUT usage, as the unification unit is shared.

\textbf{Domain Scaling}: Beyond 64-bit domains, hierarchical bit structures enable scaling to 262K values:
\begin{itemize}
  \item \textbf{Hierarchical4kChirho} (4,096 values): 64$\times$64 bit tree, 39 ns intersect
  \item \textbf{Hierarchical65kChirho} (65,536 values): 256$\times$256 bit tree (2-level), 181 ns intersect
  \item \textbf{Hierarchical256kChirho} (262,144 values): 64$\times$64$\times$64 tree (3-level), 2.5 $\mu$s intersect
  \item \textbf{Hierarchical262kWideChirho} (262,144 values): 512$\times$512 tree (2-level), \textbf{1.7 $\mu$s intersect}
\end{itemize}
The root-level sparse check eliminates subtrees early, maintaining sub-linear scaling for typical constraint patterns.

\textbf{Key finding (January 2026)}: The 2-level 512$^2$ structure is \textbf{1.5$\times$ faster} than the 3-level 64$^3$ structure for the same 262K capacity, due to fewer memory indirections:

\begin{table}[H]
\centering
\caption{Hierarchical architecture comparison for 262K-value domains}
\label{tab:hierarchical_e_chirho}
\begin{tabular}{lrrrrr}
\toprule
\textbf{Structure} & \textbf{Levels} & \textbf{10\%} & \textbf{50\%} & \textbf{90\%} & \textbf{Avg Speedup} \\
\midrule
64$^3$ (deep) & 3 & 892 ns & 2,649 ns & 4,319 ns & 1.0$\times$ \\
\textbf{512$^2$ (wide)} & \textbf{2} & \textbf{554 ns} & \textbf{1,702 ns} & \textbf{2,967 ns} & \textbf{1.5$\times$} \\
\bottomrule
\end{tabular}
\end{table}

The 512-bit words align with AVX-512 SIMD instructions, enabling single-instruction intersection on supporting CPUs. For FPGA, the wider words reduce control logic complexity at the cost of larger register files.

% ============================================================================
\section{Synthesis Results}
\label{sec:synthesis_chirho}
% ============================================================================

We present both simulation-based latency measurements and actual synthesis
results from Yosys.

\subsection{Synthesis Environment}
\label{subsec:synthesis_env_chirho}

\begin{itemize}
  \item \textbf{Synthesis Tool}: Yosys 0.61 (open-source)
  \item \textbf{Simulation Tool}: Verilator 5.x for cycle-accurate timing
  \item \textbf{Source}: Clash-generated \texttt{searchEngineChirho.v} (35 KB)
  \item \textbf{Target}: Intel Cyclone V 5CSEBA6U23I7 (DE10-Nano compatible)
\end{itemize}

\subsection{Actual Synthesis Results}
\label{subsec:actual_synthesis_chirho}

Table~\ref{tab:yosys_results_chirho} reports resource utilization from Yosys
synthesis of the complete searchEngineChirho design (January 2026).

\begin{table}[ht]
\centering
\caption{Yosys synthesis results for searchEngineChirho}
\label{tab:yosys_results_chirho}
\begin{tabular}{lr}
\toprule
\textbf{Resource} & \textbf{Count} \\
\midrule
Total Cells & 43,136 \\
Flip-Flops (DFFE) & 9,781 \\
AND Gates (ANDNOT + AND) & 16,579 \\
Multiplexers (MUX) & 4,822 \\
NOT Gates & 1,100 \\
NOR Gates & 481 \\
NAND Gates & 118 \\
Wire Bits & 127,487 \\
Port Bits & 587 \\
\bottomrule
\end{tabular}
\end{table}

\textbf{Target Device Fit}:
\begin{itemize}
  \item Intel Cyclone V (110K ALMs): $\sim$8\% utilization
  \item AMD Artix-7 A35 (33K LUTs): $\sim$65\% utilization
  \item Lattice iCE40 HX8K (8K LUTs): Does not fit
\end{itemize}

\textbf{Calyx Implementation Results.}
The Calyx-generated Verilog synthesizes after stripping simulation-only
\texttt{\$fatal} constructs:

\begin{table}[ht]
\centering
\caption{Yosys synthesis results for Calyx unify\_unit\_chirho}
\label{tab:calyx_results_chirho}
\begin{tabular}{lr}
\toprule
\textbf{Resource} & \textbf{Count} \\
\midrule
Total Cells (main + unify\_unit) & 627 \\
Flip-Flops & 203 \\
AND Gates & 140 \\
MUX & 64 \\
\bottomrule
\end{tabular}
\end{table}

The Calyx design is intentionally smaller---it implements only the unification
core, while the Clash design includes hash-consing, branching, and backtracking.

\subsection{AWS F2 Physical Synthesis (January 2026)}
\label{subsec:aws_f2_chirho}

We completed physical synthesis on AWS F2 infrastructure using Vivado 2025.1.
Table~\ref{tab:f2_results_chirho} summarizes the results.

\begin{table}[ht]
\centering
\caption{AWS F2 synthesis results (Vivado 2025.1)}
\label{tab:f2_results_chirho}
\begin{tabular}{lr}
\toprule
\textbf{Parameter} & \textbf{Value} \\
\midrule
Target FPGA & AMD VU47P-HBM \\
Clock Frequency & 250 MHz (4 ns period) \\
Timing Slack (WNS) & +0.159 ns \\
Timing Status & \textbf{MET} \\
Build Time & 50 min 47 sec \\
Instance Type & c5.9xlarge (synthesis) \\
\midrule
Initial AFI ID & afi-0710e3d924fa21ff4 \\
Initial AGFI & agfi-04abde24231f6775c \\
\textbf{Benchmark AGFI} & \textbf{agfi-05988b0b1980d6d2f} \\
AFI Status & Available \\
\bottomrule
\end{tabular}
\end{table}

\textbf{Verification on Physical Hardware.}
The AFI was loaded onto an f2.6xlarge instance and verified via the OCL
(AXI-Lite) register interface:

\begin{itemize}
  \item \texttt{VERSION (0x000)}: Read 0xF2010001 (correct design ID)
  \item \texttt{CONTROL (0x004)}: Read 0x00000000 (idle state)
  \item \texttt{STATUS (0x008)}: Read 0x00000003 (ready, bits 0-1 set)
  \item \texttt{CMD\_LO (0x010)}: Write 0xDEADBEEF, read back correctly
\end{itemize}

This confirms the design operates correctly on physical hardware at 250 MHz.

\subsection{Benchmark Methodology}
\label{subsec:methodology_chirho}

All benchmarks use \texttt{agfi-05988b0b1980d6d2f} unless otherwise noted.

\textbf{Timing Measurement.}
\begin{itemize}
  \item \textbf{Timer}: \texttt{clock\_gettime(CLOCK\_MONOTONIC)} on x86 host
  \item \textbf{Resolution}: $\sim$1 nanosecond
  \item \textbf{Scope}: First PCIe write to final PCIe read
\end{itemize}

\textbf{What Timing Includes.}
\begin{itemize}
  \item FPGA kernel execution (unification/search computation)
  \item PCIe register read/write overhead
  \item FPGA internal memory access (BRAM, HBM where applicable)
\end{itemize}

\textbf{What Timing Excludes.}
\begin{itemize}
  \item Host-side data preparation (parsing, struct packing)
  \item DMA bulk transfers (not used in register-based benchmarks)
  \item Result post-processing and file I/O
\end{itemize}

\textbf{Definition of ``Processed.''}
High throughput figures (e.g., 2.6B words/sec) measure parallel bitmask
comparison, not full NLP. Each ``word'' is a 16-byte binary record containing
\texttt{lemma\_id}, \texttt{strong\_num}, \texttt{verse\_id}, and position.
The FPGA performs integer comparison (\texttt{lemma\_id == target AND
verse\_id == same\_verse}). This is filtering, not string parsing.

\textbf{Verification Method.}
``Verified'' means CPU and FPGA return identical match counts for the same
query. For proximity searches, the first 10 matches are spot-checked for
correctness.

\subsection{Verified Benchmark Results}
\label{subsec:verified_benchmarks_chirho}

Table~\ref{tab:verified_results_chirho} presents verified results from
comprehensive benchmarks on real data.

\begin{table}[ht]
\centering
\caption{Verified benchmark results on physical F2 hardware}
\label{tab:verified_results_chirho}
\begin{tabular}{lrrr}
\toprule
\textbf{Benchmark} & \textbf{CPU Time} & \textbf{FPGA Time} & \textbf{Speedup} \\
\midrule
\multicolumn{4}{c}{\textit{TestForge: Constraint-Based Test Generation}} \\
10K records, 30 constraints & 61.5 ms & 0.0023 ms & \textbf{26,288$\times$} \\
5K records, 20 constraints & 4.99 ms & 0.0021 ms & \textbf{2,410$\times$} \\
\midrule
\multicolumn{4}{c}{\textit{ConfigGuard: Configuration Validation}} \\
5K files, 50 rules & 16.9 ms & 0.0018 ms & \textbf{9,657$\times$} \\
2K files, 25 rules & 3.37 ms & 0.0018 ms & \textbf{1,894$\times$} \\
\midrule
\multicolumn{4}{c}{\textit{Philologos: Greek NT Search (137K words)}} \\
8000-search batch & 662.38 ms & 0.099 ms & \textbf{6,679$\times$} \\
Single proximity search & 0.087 ms & 0.058 ms & 1.5$\times$ \\
\midrule
\multicolumn{4}{c}{\textit{Gradient: Real Fixed-Point Differentiable Logic$^\dagger$}} \\
50 vars, Gumbel-softmax SAT & 0.36 ms & 0.0019 ms & \textbf{185$\times$} \\
100 vars, Relaxed SAT + annealing & 0.16 ms & 0.0019 ms & \textbf{82$\times$} \\
\bottomrule
\multicolumn{4}{l}{\footnotesize $^\dagger$Q16.16 fixed-point: real gradients via chain rule, Gumbel-softmax reparameterization,} \\
\multicolumn{4}{l}{\footnotesize temperature annealing (exp/linear/cosine), scaled dot-product attention.}
\end{tabular}
\end{table}

\textbf{Key Findings.}
\begin{enumerate}
  \item \textbf{Batch mode is essential}: Single operations show 1.5$\times$
    speedup (PCIe-limited), but batching 8000 searches achieves 6,679$\times$.
  \item \textbf{Constraint complexity matters}: TestForge with FK/unique
    constraints shows highest speedup (26K$\times$) as FPGA parallelizes
    constraint checking.
  \item \textbf{Real data verified}: Greek NT proximity searches (e.g.,
    ``$\chi\rho\iota\sigma\tau o\varsigma$ within 8 words of
    $\iota\eta\sigma o\upsilon\varsigma$'') verified with CPU producing
    identical match counts.
\end{enumerate}

\textbf{Verified Proximity Search Results.}
Using the Greek New Testament (MorphGNT, 137,498 words):

\begin{table}[ht]
\centering
\caption{Verified Greek NT proximity searches}
\label{tab:greek_nt_chirho}
\begin{tabular}{lrrrr}
\toprule
\textbf{Search} & \textbf{Distance} & \textbf{Matches} & \textbf{CPU} & \textbf{FPGA} \\
\midrule
$\theta\varepsilon o\varsigma$ near $\lambda o\gamma o\varsigma$ & 8 words & 59 & 98.2 ms & 0.012 ms \\
$\chi\rho\iota\sigma\tau o\varsigma$ near $\iota\eta\sigma o\upsilon\varsigma$ & 15 words & 239 & 63.4 ms & 0.012 ms \\
$\kappa\upsilon\rho\iota o\varsigma$ near $\theta\varepsilon o\varsigma$ & 20 words & 120 & 78.8 ms & 0.012 ms \\
$\pi\iota\sigma\tau\iota\varsigma$ near $\chi\rho\iota\sigma\tau o\varsigma$ & 25 words & 40 & 61.4 ms & 0.012 ms \\
\bottomrule
\end{tabular}
\end{table}

These match counts were verified by running identical searches on CPU and
confirming agreement.

\subsection{Simulation Environment}
\label{subsec:simulation_chirho}

For cycle-accurate timing, we use:
\begin{itemize}
  \item \textbf{Tool}: Calyx 0.7.x $\rightarrow$ Verilog $\rightarrow$ Verilator 5.x
  \item \textbf{Testbench}: \texttt{calyx\_chirho/tb\_domain\_chirho.cpp}
\end{itemize}

\subsection{Latency Measurements}
\label{subsec:latency_chirho}

Table~\ref{tab:latency_chirho} reports cycle counts from simulation.

\begin{table}[ht]
\centering
\caption{Operation latency (simulation cycles)}
\label{tab:latency_chirho}
\begin{tabular}{lr}
\toprule
\textbf{Operation} & \textbf{Cycles} \\
\midrule
Domain intersection (unify) & 2 \\
Failure check (zero-test) & 1 \\
Branch (fork + stack push) & 3 \\
Backtrack (stack pop) & 2 \\
Hash cons lookup (hit) & 1 \\
Hash cons insert (miss) & 2 \\
Full unify sequence (init + compute + store) & 8 \\
\bottomrule
\end{tabular}
\end{table}

At 100 MHz, the 8-cycle unify sequence completes in 80 ns---approximately
37$\times$ faster than the software implementation's 3 $\mu$s on a 3 GHz CPU.

\subsection{Verification Results}
\label{subsec:verification_chirho}

The Verilator testbench (\texttt{tb\_domain\_chirho.cpp}) verifies:

\begin{enumerate}
  \item \texttt{0xFF AND 0xF0 = 0xF0}: Domain intersection produces correct result
  \item Failure detection: Zero result sets \texttt{failed\_chirho} flag
  \item Singleton check: $x \land (x-1) = 0$ correctly identifies single-value domains
\end{enumerate}

All tests pass in simulation. The full test completes in 8 clock cycles as
reported by the testbench.

\subsection{Performance Comparison}
\label{subsec:performance_chirho}

Table~\ref{tab:comparison_chirho} summarizes the headline performance comparison
from verified benchmarks on AWS F2 at 250 MHz.

\begin{table}[ht]
\centering
\caption{FPGA vs CPU performance for finite-domain unification (verified)}
\label{tab:comparison_chirho}
\begin{tabular}{lrrr}
\toprule
\textbf{Metric} & \textbf{FPGA (250 MHz)} & \textbf{CPU (3 GHz)} & \textbf{Speedup} \\
\midrule
Batch throughput & 40.5M unify/sec & 500K unify/sec & \textbf{81$\times$} \\
Per-op latency & 24.7 ns & $\sim$2 $\mu$s & \textbf{81$\times$} \\
8000-search batch & 0.099 ms & 662 ms & \textbf{6,679$\times$} \\
TestForge (10K rec) & 0.0023 ms & 61.5 ms & \textbf{26,288$\times$} \\
ConfigGuard (5K files) & 0.0018 ms & 16.9 ms & \textbf{9,657$\times$} \\
\bottomrule
\end{tabular}
\end{table}

\textbf{Key advantages of FPGA:}
\begin{itemize}
  \item \textbf{Deterministic latency}: No cache misses, GC pauses, or branch mispredictions
  \item \textbf{Batch efficiency}: Amortizes PCIe overhead across thousands of operations
  \item \textbf{Parallel constraint checking}: FK/unique constraints checked simultaneously
  \item \textbf{Energy efficiency}: $\sim$1 W vs $\sim$100 W for comparable throughput
\end{itemize}

\textbf{Important caveat}: Single-operation speedup is modest (1.5$\times$) due to
PCIe round-trip latency ($\sim$1.1 $\mu$s). The dramatic speedups (6,679$\times$+)
require batching multiple operations to amortize communication overhead.

% ============================================================================
\section{Production-Scale Training with HBM Batching}
\label{sec:production_scale_chirho}
% ============================================================================

The benchmarks in Section~\ref{subsec:verified_benchmarks_chirho} measure
search/inference workloads. For training differentiable models (Paper~F),
we developed an HBM-batched approach that eliminates PCIe overhead per iteration.

\subsection{The PCIe Bottleneck}
\label{subsec:pcie_bottleneck_chirho}

Naive FPGA training transfers weights and gradients via PCIe \emph{every iteration}:
\begin{itemize}
  \item PCIe Gen3 $\times$16: $\sim$15 GB/s theoretical, $\sim$1.1 $\mu$s per small transfer
  \item Training iteration (500 iters $\times$ 50 samples): 25,000 PCIe round-trips
  \item PCIe overhead dominates: $25{,}000 \times 1.1~\mu\text{s} = 27.5~\text{ms}$
\end{itemize}

The ``speedup'' of measuring PCIe transfers is misleading---it measures communication, not compute.

\subsection{HBM Batching Architecture}
\label{subsec:hbm_batch_chirho}

AWS F2's VU47P-HBM provides 8 GB HBM at 460 GB/s aggregate bandwidth.
Our HBM-batched training:

\begin{enumerate}
  \item \textbf{Upload once}: Weights, clauses, and training config $\to$ HBM
  \item \textbf{Train in HBM}: All 500 iterations $\times$ 50 samples execute in FPGA
  \item \textbf{Download once}: Final trained weights $\to$ host
\end{enumerate}

\begin{table}[ht]
\centering
\caption{PCIe transfers: Naive vs.\ HBM batched (500 iterations, 50 samples)}
\label{tab:pcie_comparison_chirho}
\begin{tabular}{lrr}
\toprule
\textbf{Approach} & \textbf{PCIe Transfers} & \textbf{Overhead} \\
\midrule
Naive (per-iteration) & 25,000 & 27.5 ms \\
HBM Batched & 2 & 2.2 $\mu$s \\
\bottomrule
\end{tabular}
\end{table}

\subsection{Production-Scale Benchmark Results}
\label{subsec:production_results_chirho}

Table~\ref{tab:production_chirho} presents verified benchmarks using HBM batching
on the production AFI (\texttt{agfi-05988b0b1980d6d2f}).

\begin{table}[ht]
\centering
\caption{Production-scale training benchmarks (HBM batched, Q16.16 fixed-point)}
\label{tab:production_chirho}
\begin{tabular}{llrrrr}
\toprule
\textbf{Mode} & \textbf{Scenario} & \textbf{Vars} & \textbf{Clauses} & \textbf{CPU (ms)} & \textbf{Speedup} \\
\midrule
\multicolumn{6}{c}{\textit{SAT Training (500 iterations, 50--200 samples)}} \\
Train & Small SAT & 100 & 300 & 200.8 & \textbf{76,633$\times$} \\
Train & Medium SAT & 500 & 1500 & 1,982 & \textbf{786,637$\times$} \\
Train & Large SAT & 1000 & 3000 & 7,916 & \textbf{3,311,925$\times$} \\
Train & XL SAT & 2000 & 6000 & 3,167 & \textbf{1,115,204$\times$} \\
\midrule
\multicolumn{6}{c}{\textit{Attention Training (50--100 iterations, 4--8 heads)}} \\
Train-Attn & Small Attn & 64 & 128 & 84.3 & \textbf{49,897$\times$} \\
Train-Attn & Medium Attn & 128 & 256 & 965 & \textbf{442,631$\times$} \\
Train-Attn & Large Attn & 256 & 512 & 3,467 & \textbf{1,469,213$\times$} \\
\midrule
\multicolumn{6}{c}{\textit{Inference (100 iterations)}} \\
Infer & Light Load & 100 & 300 & 19.1 & \textbf{11,367$\times$} \\
Infer & Medium Load & 500 & 1500 & 96.3 & \textbf{50,669$\times$} \\
Infer & Heavy Load & 1000 & 3000 & 192.2 & \textbf{107,376$\times$} \\
Infer & Burst Load & 500 & 1500 & 480.6 & \textbf{221,476$\times$} \\
\bottomrule
\end{tabular}
\end{table}

\textbf{Key insight}: The 3.3 million$\times$ speedup for large SAT training
reflects the true compute advantage when PCIe overhead is amortized via HBM batching.

\subsection{Hybrid Boolean/Differentiable Architecture}
\label{subsec:hybrid_arch_chirho}

The FPGA now hosts two complementary engines:

\begin{enumerate}
  \item \textbf{searchEngineChirho}: Fast Boolean solver (AND/OR/mask at 250 MHz)
  \item \textbf{diffTrainChirho}: Differentiable training unit (Q16.16 Gumbel-softmax)
\end{enumerate}

These work together for neural-guided SAT solving:
\begin{itemize}
  \item \textbf{Offline}: diffTrainChirho learns variable selection heuristics from problem structure
  \item \textbf{Online}: searchEngineChirho performs fast exact search using learned heuristics
  \item \textbf{Feedback}: Search failures improve heuristic training (reinforcement loop)
\end{itemize}

The DiffTrainChirho module (Clash, synthesizable) provides:
\begin{itemize}
  \item Q16.16 and Q8.8 fixed-point arithmetic via DSP blocks
  \item Taylor-series exp/log approximation (4 terms, $\pm$5\% accuracy for $|x| < 2$)
  \item LFSR Gumbel noise generation (32-bit maximal-length polynomial)
  \item Gradient accumulation with saturation arithmetic
  \item 8-state FSM: Idle $\to$ Load $\to$ Sample $\to$ Eval $\to$ Accum $\to$ Update $\to$ Store $\to$ Done
\end{itemize}

See \texttt{clash\_chirho/DiffTrainChirho.hs} for the complete implementation.

% ============================================================================
\section{Limitations and Future Work}
\label{sec:limitations_chirho}
% ============================================================================

\subsection{Current Limitations}
\label{subsec:current_limits_chirho}

\begin{enumerate}
  \item \textbf{Tabling in Software}: Mode analysis and tabling (Paper D) remain
    in software. The hardware accelerates the inner loop, but relation caching
    and goal scheduling run on the host CPU.

  \item \textbf{Timing Verified (Resolved)}: As of January 2026, physical synthesis
    with Vivado 2025.1 achieved timing closure at 250 MHz on AWS F2 (AMD VU47P-HBM).
    The AFI was successfully loaded and verified on physical hardware.

  \item \textbf{Fixed Variable Count}: The current design supports 8 variables.
    Scaling to hundreds requires architectural changes (memory banking, pipelining).

  \item \textbf{No Occurs Check}: The hardware unification unit does not implement
    the occurs check. Cyclic structures may cause non-termination.

  \item \textbf{Sequential Constraint Processing}: Constraints are processed
    one at a time. Parallel constraint propagation would require multiple
    unification units with conflict detection.
\end{enumerate}

\subsection{Future Work}
\label{subsec:future_chirho}

\begin{enumerate}
  \item \textbf{Physical Synthesis}: \emph{Completed.} AWS F2 synthesis achieved
    timing closure at 250 MHz (Section~\ref{subsec:aws_f2_chirho}).

  \item \textbf{Hardware Tabling}: Move relation caching to BRAM with
    hardware-managed LRU eviction.

  \item \textbf{Parallel Search}: Instantiate multiple search engines sharing
    the hash cons unit for embarrassingly parallel exploration.

  \item \textbf{ASIC Exploration}: The highly regular structure suggests
    potential for ASIC implementation with sub-nanosecond unification.
\end{enumerate}

% ============================================================================
\section{Related Work}
\label{sec:related_chirho}
% ============================================================================

\textbf{Prolog on Hardware}: Early work on Prolog machines (PSI, PIM) used
specialized architectures for unification. Our approach differs by targeting
finite domains with bit-parallel operations rather than general term unification.

\textbf{SAT/SMT Accelerators}: FPGA SAT solvers (e.g., Catapult, SatHardware)
accelerate Boolean satisfiability. Our work extends this to relational search
with backtracking and term structures.

\textbf{Constraint Acceleration}: GPU-accelerated constraint propagation
(OpenCL, CUDA) achieves high throughput but lacks deterministic latency.
FPGAs offer predictable timing critical for real-time applications.

\textbf{Calyx and Clash}: We leverage modern hardware compilation infrastructure:
Calyx~\cite{calyx} for correct-by-construction lowering and Clash~\cite{clash}
for verified Haskell-to-Verilog translation.

% ============================================================================
\section{Conclusion}
\label{sec:conclusion_chirho}
% ============================================================================

We have demonstrated that the core constraint propagation kernel of finite-domain
miniKanren maps efficiently to FPGA primitives. Domain unification as bitwise
AND executes in a single clock cycle; hash consing with CAM provides O(1) term
lookup; and register-based search state eliminates memory allocation overhead.

As of January 2026, we have validated these results on physical AWS F2 hardware
(AMD VU47P-HBM) at 250 MHz. Comprehensive benchmarks on real data---including
the Greek New Testament (137,498 words)---demonstrate:

\begin{itemize}
  \item \textbf{Batch throughput}: 40.5M unifications/second (6,679$\times$ vs CPU)
  \item \textbf{TestForge}: 26,288$\times$ speedup for 10K-record constraint solving
  \item \textbf{ConfigGuard}: 9,657$\times$ speedup for 5K-file validation
  \item \textbf{Philologos}: Verified proximity searches on Greek NT with CPU agreement
\end{itemize}

\textbf{Key insight}: Single-operation speedup is modest (1.5$\times$) due to PCIe
latency. The dramatic speedups require \emph{batch mode}, where thousands of
operations amortize communication overhead. This is the natural operating mode
for SaaS applications serving concurrent users.

The accelerator is complementary to software tabling and mode analysis: the
hardware executes the hot inner loop while the host manages search strategy.
This hybrid approach preserves the expressive power of relational programming
while enabling hardware-level performance for production workloads.

\paragraph{Artifacts.}
\begin{itemize}
  \item Calyx implementation: \texttt{calyx\_chirho/domain\_chirho.futil}
  \item Clash Boolean solver: \texttt{clash\_chirho/MiniKanrenChirho.hs}, \texttt{clash\_chirho/HashConsChirho.hs}
  \item \textbf{Clash differentiable training}: \texttt{clash\_chirho/DiffTrainChirho.hs}
  \item Verilator testbench: \texttt{calyx\_chirho/tb\_domain\_chirho.cpp}
  \item AWS F2 CL wrapper: \texttt{synth\_chirho/aws\_f2\_chirho\_cl/}
  \item Timing report: \texttt{synth\_chirho/aws\_f2\_chirho\_cl/build\_results\_chirho/}
  \item Initial AFI: \texttt{agfi-04abde24231f6775c} (validation)
  \item \textbf{Benchmark AFI}: \texttt{agfi-05988b0b1980d6d2f} (comprehensive benchmarks)
\end{itemize}

\paragraph{Benchmark Data Files.}
All raw data in \texttt{synth\_chirho/aws\_f2\_chirho\_cl/benchmarks\_chirho/}:
\begin{itemize}
  \item \texttt{benchmark\_results\_chirho.csv} --- Core FPGA register/HBM tests
  \item \texttt{comprehensive\_saas\_results\_chirho.csv} --- 105 SaaS scenarios (FPGA projected)
  \item \texttt{comprehensive\_verified\_chirho.csv} --- \textbf{60 scenarios verified CPU vs FPGA}
  \item \texttt{philologos\_verified\_results\_chirho.csv} --- Greek NT searches (verified)
  \item \texttt{batch\_8000\_results\_chirho.csv} --- 8000-search batch test
  \item \texttt{saas\_verified\_results\_chirho.csv} --- CPU vs FPGA verification
  \item \texttt{COMPREHENSIVE\_BENCHMARK\_REPORT\_CHIRHO.md} --- Full methodology and results
  \item \textbf{\texttt{production\_scale\_results\_chirho.csv}} --- Production training (3.3M$\times$ speedup)
  \item \textbf{\texttt{gradient\_verified\_chirho.csv}} --- Differentiable logic benchmarks
\end{itemize}

\paragraph{Code Availability.}
\url{https://github.com/loveJesus/minikanren_1bit_chirho}

\paragraph{Companion Papers.}
This paper is part of a suite:
\begin{itemize}
  \item \paperA
  \item \paperB
  \item \paperC
  \item \paperD
  \item \paperF
\end{itemize}

% ============================================================================
\section*{Acknowledgments}
% ============================================================================

\standardacknowledgments

I also thank the Calyx and Clash communities for their excellent tools and
documentation.

% ============================================================================
% References
% ============================================================================

\begin{thebibliography}{9}

\bibitem{calyx}
Rachit Nigam, Samuel Thomas, Zhijing Li, and Adrian Sampson.
\newblock Calyx: A Language for Hardware Accelerator Design.
\newblock In \emph{PLDI}, 2021.

\bibitem{clash}
Christiaan Baaij.
\newblock Digital Circuit Design in Haskell.
\newblock PhD thesis, University of Twente, 2015.

\bibitem{reasoned}
Daniel P. Friedman, William E. Byrd, and Oleg Kiselyov.
\newblock \emph{The Reasoned Schemer}.
\newblock MIT Press, 2005.

\bibitem{artix7}
Xilinx.
\newblock 7 Series FPGAs Data Sheet: Overview (DS180).
\newblock 2020.

\end{thebibliography}

\vspace{1em}
\begin{center}
\textit{Soli Deo Gloria} $\chi\rho$
\end{center}

\end{document}
